\section{The infinite-height coordinate predual}\label{binf:definition}

The finite-height forest and the infinite tree have the same finite
segment tests, but their completions must be specified separately.
Here we give the infinite-tree construction and the precise finite-head
facts needed by the subsequent arguments.

Let $T_\infty=\mathbb N^{<\omega}$, with root $o$, ordered by extension
of finite words. On $c_{00}(T_\infty)$ use the real segment norm
\[
 \rho(u)=\sup_{\mathcal S}
 \left(\sum_{S\in\mathcal S}u(S)^2\right)^{1/2},
 \qquad u(S)=\sum_{v\in S}u(v),
\]
where $\mathcal S$ ranges over finite families of pairwise disjoint
nonempty segments. Singleton tests and the Euclidean triangle inequality show that this is
a norm, with $\|u\|_2\le\rho(u)\le\|u\|_1$.
Let $J_\infty$ be its completion and define
\[
 B_\infty=\overline{\operatorname{span}}\{e_v^*:v\in T_\infty\}
 \subset J_\infty^*.
\]
The closure and norm in this definition are those of $J_\infty^*$.
This is the coordinate predual described in
\cite[Section~2.II]{Banakh2000}. On finite coefficient vectors its norm is
\begin{equation}\label{binf:polar}
 \left\|\sum_v a_v e_v^*\right\|
 =\sup\left\{\left|\sum_v a_v u(v)\right|:
                 u\in c_{00}(T_\infty),\ \rho(u)\le1\right\}.
\end{equation}
Indeed finite tests are dense in $J_\infty$, with rescaling permitting
approximation within its unit ball. Consequently completing the finite
coefficient vectors in \eqref{binf:polar} gives exactly $B_\infty$.
In particular $\|a\|_\infty\le\|a\|\le\|a\|_2$ for finite
coefficient arrays, and every coordinate functional has norm one, by
testing on the corresponding primal coordinate vector.
For any $x\in B_\infty$, the same supremum computes $\|x\|$:
each test is continuous, and the supremum is a $1$-Lipschitz function
of $x$ agreeing with its norm on the dense coordinate span. We use only
this closed coordinate span, without identifying it with the full dual
$J_\infty^*$.

\begin{lemma}[Finite heads and their complements]\label{binf:projections}
Let $H\subset T_\infty$ be finite and ancestral. The coordinate
projections $P_H$ and $Q_H=I-P_H$ are contractive on $J_\infty$ and
$B_\infty$. The range of $P_H$ is finite dimensional, and finite
coefficient vectors supported outside $H$ are dense in $\ker P_H$.
Along every increasing ancestral exhaustion, $P_Hx\to x$ in
$B_\infty$. A nonzero vector $x=P_Hx$ has a norming test $f=P_Hf$
with $\rho(f)=1$ and $x(f)=\|x\|$.
\end{lemma}
\begin{proof}
A segment intersects $H$ in a prefix and its complement in a suffix.
Restricting every member of a disjoint testing family proves the primal
contractions. Their adjoints preserve the coordinate span, hence its
closure, and give the asserted dual contractions. Applying $Q_H$ to
finite approximants proves density in its range. An exhaustion eventually
fixes each finite coefficient vector; contractivity extends convergence
to the completion. Finally the norm of $x=P_Hx$ can be computed using
only $P_HJ_\infty$, whose unit ball is compact. A maximizer is a unit
test when $x\ne0$.
\end{proof}

For a segment $S\subset T_\infty$ write $d_S=\sum_{v\in S}e_v^*$.
A represented atom is a finite sum $\sum_i a_i d_{S_i}$ with disjoint
segments $S_i$ and $\sum_i a_i^2\le1$. Cauchy--Schwarz gives its norm
at most one. The representation is part of the data when coefficients
or crossings are subsequently examined.

\begin{lemma}[Finite polar hull and completed atomic ball]
\label{binf:finite-hull}
For finite ancestral $H$, the unit ball of $P_HB_\infty$ is the convex
hull of represented atoms with all segments in $H$. Each point is a
convex combination of at most $|H|+1$ such atoms. The full unit ball
of $B_\infty$ is the norm-closed convex hull of finite atoms.

On $P_HB_\infty$ the norm also equals the gauge
\[
 \inf\left\{\sum_j\lambda_j:
   x=\sum_j\lambda_j A_j,\quad \lambda_j\ge0,\quad
   A_j\text{ atoms supported in }H\right\}.
\]
In particular, a vector of norm strictly below one has such a finite
representation with $\sum_j\lambda_j<1$. The gauge defined on
$c_{00}(T_\infty)$ by allowing all finite atoms has the same value.
Its completion is therefore isometric to $B_\infty$.
\end{lemma}
\begin{proof}
There are finitely many disjoint segment families in $H$. For each
family, its atoms form the continuous image of a Euclidean unit ball.
Their union is compact, as is its finite-dimensional convex hull $C$.
For a test $f=P_Hf$, the support function of $C$ is precisely $\rho(f)$,
by optimizing the coefficients in each family. Finite-dimensional
separation identifies $C$ with the dual unit ball of $P_HJ_\infty$;
contractivity identifies this with the inherited ball of $P_HB_\infty$.
Carath\'eodory's theorem gives $|H|+1$ atoms. Equivalently, its elementary
proof repeatedly removes a term from any longer convex combination
using an affine dependence among its points, keeping the coefficients
nonnegative and their sum one.

For the gauge assertion, every displayed representation bounds the
norm by $\sum_j\lambda_j$. Conversely, write $x=r(x/r)$, where
$r=\|x\|>0$, and represent $x/r\in C$; $x=0$ is immediate.
For the global gauge, choose an ancestral $H$ containing the support
of $x$. Projection $P_H$ sends every global atom to an atom in $H$:
intersect its disjoint segments with $H$ and discard empty intersections.
Thus every global representation projects to a local representation
with no larger total weight. Local representations are already global
ones, so the two gauges agree on every finite coefficient vector.
The completion identification follows from \eqref{binf:polar}.

For a completed unit vector, Lemma~\ref{binf:projections} gives
$P_Hx\to x$ with each $P_Hx$ in the finite atomic ball. Thus the
norm-closed convex hull of all finite atoms is the completed unit ball.
This last assertion is an approximation statement; the finite exact
representations are used only after a finite projection.
\end{proof}

Let $G_H$ be the minimal vertices outside $H$, and let
$U_g=\{v:g\preceq v\}$ for $g\in G_H$. When $H$ is empty the
root is the sole gate.

\begin{lemma}[Exterior square sums]\label{binf:cones}
For a finite test $u=Q_Hu$,
\begin{equation}\label{binf:primal-cones}
 \rho(u)^2=\sum_{g\in G_H}\rho(u\mathbf1_{U_g})^2.
\end{equation}
After completion,
$Q_HJ_\infty=(\bigoplus_{g\in G_H}J_\infty(U_g))_{\ell_2}$.
For $x\in B_\infty$,
\begin{equation}\label{binf:dual-cones}
 \|Q_Hx\|^2=\sum_{g\in G_H}\|x|_{J_\infty(U_g)}\|^2.
\end{equation}
The map $x\mapsto(x(e_g))_{g\in G_H}$ is a contraction to
$\ell_2(G_H)$.
\end{lemma}
\begin{proof}
Each exterior part of a segment lies in one cone. Grouping tests by
cones gives one inequality in \eqref{binf:primal-cones}; joining nearly
maximizing disjoint families from the finitely many nonzero cones gives
the other. Completion gives the displayed primal sum. The norm of
$Q_Hx$ is exactly that of its restriction to $Q_HJ_\infty$, by
contractivity. The dual norm of an $\ell_2$ sum is the $\ell_2$ sum
of the restriction norms: Cauchy--Schwarz proves the upper bound, and
finitely many almost norming tests with proportional weights give the
lower bound. This proves \eqref{binf:dual-cones} without any assertion
that $B_\infty$ is the full dual. Finally a segment meets the antichain
$G_H$ in at most one point, so tests supported there have precisely their
Euclidean norm. Testing $x$ on them gives the contraction. The same
argument shows that $x\mapsto(x(e_v))_{v\in A}$ is a contraction
to $\ell_2(A)$ for every antichain $A$, not only for a gate set.
\end{proof}

